% !PRINTED 548
\weberpage{548}
%<<CONT
worin $p_1, p_2 \ldots$ Potenzen von \gesperrt{verschiedenen} Primzahlen sind.
Wir setzen noch
\[
m = p_1 m_1 = p_2 m_2 \ldots,
\]
und wählen, was nach dem soeben Bewiesenen stets möglich ist,
$\lambda_1$ relativ prim zu $p_1$, $\lambda_2$ relativ prim zu $p_2$ etc., so dass
\[
(\varepsilon^{\lambda_1}, \alpha), \quad (\varepsilon^{\lambda_2}, \alpha) \ldots
\]
von Null verschieden sind. Dann ist nach dem Theorem \S.\,164, 3.:
\weq{5}{(\varepsilon^{\lambda}, \alpha)(\varepsilon^{\lambda_1}, \alpha)^{m_1 \nu}(\varepsilon^{\lambda_2}, \alpha)^{m_2 \nu} \ldots = \chi_{\lambda}}
eine in $\Omega(\varepsilon)$ enthaltene Grösse, wenn
\weq{6}{\lambda \equiv -\nu(\lambda_1 m_1 + \lambda_2 m_2 + \cdots) \ (\mathrm{mod}\ m).}
Es ist aber $(\varepsilon^{\lambda_1}, \alpha)^{m_1}$ eine Wurzel $p_1^{\mathrm{ten}}$ Grades einer Function
in $\Omega(\varepsilon)$, und wir setzen also
% [sic] authority omits the exponent m_2 from the second factor; symmetry with the first factor requires it.
\weq{7}{(\varepsilon^{\lambda_1}, \alpha)^{m_1} = \sqrt[p_1]{\varphi_1}, \quad (\varepsilon^{\lambda_2}, \alpha) = \sqrt[p_2]{\varphi_2} \ldots, \quad (\varepsilon^{\lambda}, \alpha) = \sqrt[m]{\psi_{\lambda}}.}
Dann wird nach (5)
\weq{8}{\sqrt[m]{\psi_{\lambda}} = \frac{\chi_{\lambda}}{\left(\sqrt[p_1]{\varphi_1}\sqrt[p_2]{\varphi_2} \ldots \right)^{\nu}}.}
Nun ist $(\lambda_1 m_1 + \lambda_2 m_2 + \cdots)$ relativ prim zu $m$, da $m_2 \ldots$
durch $p_1$ theilbar, $\lambda_1 m_1$ zu $p_1$ relativ prim ist, und also erhält
man aus (6) für jedes $\lambda$ eine nach dem Modul $m$ völlig bestimmte
Zahl $\nu$.
Wenn wir also die Ausdrücke (8) in (3) einsetzen und den
Radicalen $\sqrt[p_1]{\varphi_1}, \sqrt[p_2]{\varphi_2} \ldots$ alle ihre Werthe beilegen, so erhalten
wir genau $m$ verschiedene Werthe und nicht mehr für $\alpha$.

Die letzten Resultate können wir benutzen, um eine Form
der Darstellung der Wurzeln $\alpha$ in etwas verallgemeinerter Gestalt
abzuleiten, die \gesperrt{Abel} an der angeführten Stelle mittheilt, und die
sich auf den Fall bezieht, wo der Körper $\Omega$ reell ist, d. h. aus
lauter reellen Zahlen besteht. Die Functionen $\varphi, \psi, \chi$, wie wir
sie oben benutzt haben, sind dann zusammengesetzt aus reellen
Zahlen und aus der Einheitswurzel $\varepsilon$, die man durch die Theilung der Kreisperipherie in $m$ gleiche Theile findet; man kann etwa
\[
\varepsilon = e^{\frac{2\pi i}{m}} = \cos\frac{2\pi}{m} + i \sin\frac{2\pi}{m}
\]
setzen.

Die Function $\varphi_1$ geht, wenn $\varepsilon$ in $\varepsilon^{-1}$ verwandelt wird, in
den conjugirt imaginären Werth über, den wir mit $\varphi_1'$ bezeichnen.
